\chapter{Ekosistem dan Orkestrasi Mitra}

\section*{Tujuan Pembelajaran}
\begin{itemize}
  \item Memetakan pelaku ekosistem, peran, dan aliran nilai di sekitar organisasi.
  \item Membandingkan opsi strategi platform, yaitu membangun, bergabung, atau bermitra, beserta \textit{trade-off} masing-masing.
  \item Merumuskan aturan dasar keterbukaan platform bagi mitra.
\end{itemize}

\section{Pendahuluan}

% TODO
% Cakupan: strategi "Satu Pasien, Satu Rekam Medis" bergantung pada klinik mitra, asuransi,
% laboratorium, apotek daring, dan perusahaan klien; Medikara harus memilih peran di ekosistem itu.

\section{Ekosistem sebagai Struktur Kerja Sama}

% TODO
% Cakupan: definisi ekosistem menurut Jacobides, Cennamo, dan Gawer; komplementaritas unik dan
% supermodular; modularitas sebagai syarat ekosistem; perbedaan ekosistem dengan rantai pasok dan
% aliansi.
% Kasus Medikara: hubungan saling melengkapi antara rumah sakit, laboratorium, apotek, dan asuransi.

\section{Peran dalam Ekosistem: Orkestrator, Pelengkap, dan Pengguna}

% TODO
% Cakupan: peran orkestrator, pelengkap (complementor), dan pengguna; aliran nilai berupa uang,
% data, layanan, dan reputasi; penciptaan nilai bersama dan penangkapan nilai.
% Kasus Medikara: calon peran Medikara sebagai orkestrator layanan kesehatan karyawan perusahaan.

\section{Boundary Resources: Jembatan antara Platform dan Mitra}

% TODO
% Cakupan: model boundary resources (Ghazawneh dan Henfridsson); API, SDK, dokumentasi, ketentuan
% penggunaan, dan program mitra; dua fungsi boundary resources, yaitu membuka akses dan menjaga
% kendali.
% Kasus Medikara: boundary resources yang belum dimiliki, misalnya dokumentasi API dan sandbox.
% Pendalaman teknis: menelaah portal pengembang satu platform publik dan mengelompokkan isinya.

\section{Keterbukaan dan Kendali}

% TODO
% Cakupan: platform terbuka dan tertutup; manfaat inovasi pihak luar (Parker, Van Alstyne, dan
% Jiang); risiko kehilangan kendali mutu dan keamanan; tingkat keterbukaan yang bertahap.
% Kasus Medikara: kebocoran data lewat vendor pada 2025 sebagai pelajaran tentang kendali.

\section{Tata Kelola Ekosistem}

% TODO
% Cakupan: siapa boleh bergabung; aturan perilaku mitra; pembagian nilai dan data; penyelesaian
% sengketa; kurasi mitra; tata kelola API teknis dibahas di Bab 13.
% Kasus Medikara: syarat minimum bagi klinik mitra yang ingin terhubung.

\section{Opsi Strategi: Membangun, Bergabung, atau Bermitra}

% TODO
% Cakupan: tiga opsi strategi beserta trade-off; risiko platform envelopment (Eisenmann, Parker,
% dan Van Alstyne); matriks keputusan; format catatan keputusan beserta justifikasi.
% Kasus Medikara: membangun platform sendiri dibandingkan bergabung ke platform asuransi atau
% aplikasi kesehatan nasional.

\section{Aktivitas: Peta Ekosistem dan Keputusan Strategi Platform}

% TODO
% Tujuan aktivitas: menetapkan peran organisasi dalam ekosistem dan opsi strategi platform beserta
% alasannya.
% Peran dalam tim: anggota berlatar bisnis dan hukum memimpin analisis peran dan kepentingan tiap
% pelaku; anggota berlatar teknis menandai boundary resources yang sudah ada dan yang belum ada.
% Langkah: daftar pelaku ekosistem; gambar peta ekosistem dengan aliran nilai; nilai tiga opsi
% strategi dalam matriks keputusan; tulis catatan keputusan satu halaman.
% Luaran dan tonggak capstone: peta ekosistem mitra dan catatan keputusan strategi platform,
% melengkapi C1.
% Berkas sources yang direncanakan: sources/chapter03/peta-ekosistem.md,
% sources/chapter03/matriks-opsi-strategi.md, sources/chapter03/catatan-keputusan-strategi.md

\section{Ringkasan}

% TODO

\section{Latihan}

\begin{enumerate}
  \item Jelaskan perbedaan mitra dan pelanggan dalam ekosistem layanan kesehatan, dengan satu contoh untuk masing-masing.
  \item Misalkan sebuah perusahaan asuransi mengajak Medikara bergabung ke platform milik asuransi tersebut. Analisis untung dan ruginya bagi Medikara dibandingkan membangun platform sendiri.
  \item Tentukan tiga \textit{boundary resources} yang perlu disediakan agar klinik mitra dapat terhubung tanpa pendampingan intensif.
\end{enumerate}
